\section{Introduction}
\label{sec:introduction}

Synthetic data is increasingly used to supplement scarce observations, test
data-intensive workflows, support model development, and reduce dependence on
direct access to sensitive or restricted records. In tabular learning, the
design space now spans conditional adversarial generators
\citep{xu2019ctgan}, language-model generators
\citep{borisov2023great,solatorio2023realtabformer}, and diffusion-based
approaches \citep{kotelnikov2023tabddpm,zhang2024tabsyn}. Language models also
make class-conditioned generation and task-specific prompting readily
accessible, including in imbalanced-data settings \citep{kim2024epic}. These
capabilities make generation easier to deploy, but they do not by themselves
establish that the resulting collection is fit for scientific or predictive
use.

A plausible-looking row is only one part of a usable synthetic dataset. A
candidate can violate a schema, omit a required field, leave the observed
categorical domain, contradict its requested class, or be statistically or
semantically implausible. Even records that pass such checks may create a weak
dataset through exact or near redundancy, narrow support, distorted
class-conditional structure, or poor downstream utility. Some deployments add
privacy- and fairness-specific obligations that cannot be inferred from
realism alone; this draft treats corresponding evaluators as domain-dependent
extensions rather than part of the initial empirical core. Work on constrained
tabular generation shows why distributional fit and declared domain
constraints are distinct requirements
\citep{stoian2024realistic}, while multidimensional auditing separates
fidelity, diversity, and generalization rather than treating quality as a
single number \citep{alaa2022faithful}.
% TODO[CITATION]: add primary privacy and fairness evaluation sources before
% making claims about any particular protection or fairness criterion.

Existing generation and evaluation methods illuminate different parts of this
problem. Controlled generator benchmarks hold experimental factors fixed and
combine intrinsic and downstream views \citep{kim2025evaluating}; synthetic
data surveys likewise organize generation, curation, and evaluation as
separate stages \citep{long2024survey}. Other work develops distributional
metrics, constraint enforcement, or model-based critics. These strands are
complementary, but their outputs do not share a natural scale or scope. A
schema failure is a field- or record-level fact, a judge score is an uncertain
assessment, a duplicate count is a dataset property, and downstream utility is
an outcome of a learning pipeline. Treating these signals as interchangeable
can hide both missing evidence and disagreement.

This motivates a central distinction: record validity and dataset value are
related but not equivalent. An exact copy of a valid training example may be a
well-formed member of the target class, yet contribute no new support to an
accepted synthetic collection. Conversely, a novel candidate may improve
coverage while carrying a local defect that should not be accepted unchanged.
Evaluation therefore has to ask both whether a candidate is acceptable in
isolation and what changes when that candidate is added to the dataset. The
latter question is conditional on the evolving collection and cannot be
answered by a fixed per-record score.

We study \rcqasyn{} as a sequential dataset-construction framework. A
generator proposes a candidate under a dataset specification and a
TRAIN-derived context. Record- and dataset-level evaluators emit scoped
evidence; a future calibration and fusion stage preserves uncertainty and
provenance; and a risk-aware controller makes acceptance decisions without
reducing all evidence to an undifferentiated quality score. The current
repository implements the reproducible generation, paired-baseline,
diagnostic, and human-audit infrastructure needed to study this core. It does
not yet implement the proposed M2+ evaluator suite, calibration, evidence
fusion, or controller. Localized repair and richer dataset optimization are
later extensions whose inclusion in the final system depends on evidence from
the core experiments.

\paragraph{Research questions.}
RQ1 and RQ2 are the immediate questions for the human-audited development
study. RQ3 tests the core proposed decision method after M2 implementation;
RQ4 and RQ5 are longer-range questions that need not be answered by the first
complete paper.

\begin{description}
  \item[RQ1.] To what extent do structural rules, statistical-support
    evidence, and language-model critics identify human-judged defects in
    synthetic records?
  \item[RQ2.] How does record-level validity differ from dataset-level quality,
    including redundancy, support coverage, and distributional fidelity?
  \item[RQ3.] Can calibrated heterogeneous evidence provide more reliable
    acceptance decisions than raw generation, deterministic filtering, or one
    uncalibrated language-model judge? \emph{This core question requires M2+
    implementation and final evaluation.}
  \item[RQ4.] Can defect-localized repair improve accepted-data quality more
    efficiently than complete regeneration? \emph{This extension is not yet
    implemented.}
  \item[RQ5.] Can dataset-level marginal-contribution control improve
    diversity, coverage, and downstream utility without sacrificing record
    validity? \emph{This longer-range control question remains pending.}
\end{description}

\paragraph{Intended contributions and current status.}
The current work establishes a formulation that separates record-level
validity from dataset-level quality, a reproducible paired-baseline evaluation
framework, copying and support diagnostics, and a blinded human-ground-truth
protocol. We implement these elements and evaluate the diagnostics only in
development. We propose, but have not yet validated, calibrated heterogeneous
evidence fusion and a risk-aware decision controller. Localized repair and
sophisticated marginal-contribution optimization are secondary extensions.
Accordingly, ``Risk-Calibrated'' names the intended final method and research
hypothesis; the current implementation provides no calibrated risk guarantee.
Superiority, statistical significance, and final-benchmark claims remain
open.
